\section{Background and Notation}
  \label{sec:background-and-notation}

  \subsection{The discrete Heisenberg group and its Cayley graph}

    The discrete Heisenberg group over $\mathbb{Z}/q\mathbb{Z}$ is
    \begin{equation}
      G_q = \mathrm{Heis}_3(\mathbb{Z}/q\mathbb{Z})
      = \left\{\,\begin{pmatrix}
                   1&a&\gamma\\0&1&b\\0&0&1
      \end{pmatrix}
          : a,b,\gamma\in\mathbb{Z}/q\mathbb{Z}\right\},
    \end{equation}
    with multiplication law $(a,b,\gamma)\cdot(a',b',\gamma') = (a+a',\,b+b',\,\gamma+\gamma'+ab')$.
    The group has order $|G_q|=q^3$ and centre $Z(G_q)=\{(0,0,\gamma)\}$.
    The commutator of the standard generators $X=(1,0,0)$ and $Y=(0,1,0)$ is
    $[X,Y] = (0,0,1) =: Z$, which generates the centre.

    We equip $G_q$ with the standard symmetric generating set
    $S = \{\pm X,\pm Y\}$,
    so $|S|=4$, and form the Cayley graph $\mathrm{Cay}(G_q,S)$.
    The ball $B_n = \{g\in G_q : d(e,g)\leq n\}$ satisfies
    $|B_n|\sim c\,n^4$ as $n\to\infty$ for a constant $c>0$, in accordance with
    the Bass--Guivarc'h theorem for nilpotent groups of homogeneous dimension
    $D_{\mathrm{hom}}=4$.

  \subsection{The three-step projective fingerprint}

    Let $k=3$.
    A \emph{$k$-step path} from $g\in G_q$ is a sequence $\pi=(s_1,s_2,s_3)\in S^3$
    of generators; the endpoint of $\pi$ at $g$ is $g\cdot s_1\cdot s_2\cdot s_3$.
    For a shell $S_n = B_n\setminus B_{n-1}$, the \emph{fingerprint vector} of
    a path $\pi$ starting at $g\in S_n$ is the ordered triple of endpoints
    \begin{equation}
      v(\pi,g) = (g\cdot s_1,\; g\cdot s_1\cdot s_2,\; g\cdot s_1\cdot s_2\cdot s_3)
      \;\in\; G_q^3.
    \end{equation}
    The projective capacity $\Sigma_n$ measures the effective dimension (rank) of the span
    generated by the fingerprint vectors of all paths starting in shell $S_n$,
    normalised by $|S_n|$, with $\Sigma_n\to 0$ as the span saturates the ambient space and
    new shells contribute no genuinely new directions.
    For the $k=3$ construction used throughout this paper, each element of $S_n$ contributes
    up to $|S|^3=64$ candidate directions, so $\Sigma_n$ is bounded by~$1$ only once
    $|S_n|$ exceeds the ambient dimension, not at every depth (Section~\ref{sec:exact-capacity}).
    The state law $R_n^{(k)}\approx\Phi(\bar{\eta}_n)$, with
    $\Phi(\eta)=1/\sqrt{1+\eta^2}$ and $\bar{\eta}_n=\Sigma_n/\Sigma_1$, was
    established in O9--O11~\cite{Beau2026a13,Beau2026a14,Beau2026a15}.

  \subsection{The Weil representation: exact structure vs proxy}
    \label{sec:weil-exact-vs-proxy}

    The affine representation $\rho_{\mathrm{aff}}$ of $G_q$ on $L^2(\mathbb{Z}/q\mathbb{Z})$
    decomposes under the centre $Z(G_q)\cong\mathbb{Z}/q\mathbb{Z}$ as
    \begin{equation}
      \rho_{\mathrm{aff}} = \rho_0 \oplus \bigoplus_{c=1}^{q-1} \rho_c,
    \end{equation}
    where $\rho_0$ is the trivial representation and each $\rho_c$ ($c\neq 0$) is an
    irreducible $q$-dimensional Weil representation of $G_q$, defined by~\eqref{eq:weil-exact}.
    The three-fold tensor product then splits into $q^3$ blocks:
    \begin{equation}
      \rho_{\mathrm{aff}}^{\otimes 3}
      \;\cong\; \bigoplus_{(c_1,c_2,c_3)\in(\mathbb{Z}/q\mathbb{Z})^3}
      H_{c_1,c_2,c_3},
    \end{equation}
    where each block $H_{c_1,c_2,c_3}$ has dimension $O(q)$.

    \begin{definition}[Generic blocks and notation]
      \label{def:generic}
      A multi-index $(c_1,c_2,c_3)\in(\mathbb{Z}/q\mathbb{Z})^3$ is \emph{generic} if
      $c_i\neq 0$ for $i=1,2,3$ and $c_1+c_2+c_3\not\equiv 0\pmod{q}$.
      We denote by $H_c$ any generic block $H_{c_1,c_2,c_3}$, where $c$ stands
      for the corresponding generic multi-index $(c_1,c_2,c_3)$.
      When we refer to a single central character~$c$ (for instance in the definition
      of the individual Weil representation~$\rho_c$), this is made explicit by context.
      No ambiguity arises between the two uses because $\rho_c$ refers to an individual
      irreducible factor of $\rho_{\mathrm{aff}}$, while $H_c$ always refers to an
      irreducible block of the triple tensor product.
    \end{definition}

    \paragraph{Exact Weil projection.}
      In the exact projection, each group element $(a,b,\gamma)\in G_q$ acts on
      $f\in L^2(\mathbb{Z}/q\mathbb{Z})$ via~\eqref{eq:weil-exact}, retaining the
      full phase $\psi_c(\gamma+bx)$.
      The central phase $\psi_c(\gamma)$ depends on~$\gamma$ and introduces
      a multiplicative twist that couples the central coordinate to the translational
      and phase dynamics.

    \paragraph{Bidimensional proxy of O11.}
      The proxy introduced in O11~\cite{Beau2026a15} acts on the abelianised
      coordinates $(a,b)\in(\mathbb{Z}/q\mathbb{Z})^2$ of $G_q$ via
    \begin{equation}
      (\pi_{(\tau,\sigma)}\,f)(x) = e^{2\pi i\,(\tau\cdot a + \sigma\cdot b)/q}\,f(x),
      \label{eq:proxy}
      \end{equation}
      indexed by sextuples $(\tau,\sigma)=(t_1,t_2,t_3,s_1,s_2,s_3)\in(\mathbb{Z}/q\mathbb{Z})^6$.
      The proxy ambient dimension is~$q^2$, corresponding to the $q^2$ possible pairs
      $(u_a,u_b)\in(\mathbb{Z}/q\mathbb{Z})^2$.
      This construction discards the central coordinate~$\gamma$ entirely:
      the phase $\psi_c(\gamma)$ does not appear in~\eqref{eq:proxy},
      so the proxy measures capacity in the abelianisation of the group action.

    \paragraph{Key structural distinction.}
      The difference between the exact projection and the proxy is therefore not
      merely quantitative but structural:
      the exact projection operates in a genuine irreducible representation of $G_q$,
      which is in general sensitive to the non-commutativity of the group,
      while the proxy operates in a family of one-dimensional characters of the
      abelianised group $G_q^{\mathrm{ab}}\cong(\mathbb{Z}/q\mathbb{Z})^2$.
      Whether the specific rank-based capacity observable $\Sigma_n^{(c)}$ built from the
      exact representation is \emph{itself} sensitive to this non-commutativity — as opposed
      to the representation merely containing it — is a non-trivial question that
      Section~\ref{sec:exact-capacity} answers analytically, not experimentally: it is not
      (Remark~\ref{rem:blind}).

\subsection{Block-wise projective capacity: recap of O11}
  \label{sec:o11-recap}

  O11~\cite{Beau2026a15} defined the incremental block capacity for the proxy as
  \begin{equation}
    \Sigma_n^{(\tau,\sigma)}
    = \frac{\Delta r_n^{(\tau,\sigma)}}{|S_n|},
    \qquad
    \Delta r_n^{(\tau,\sigma)}
    = \dim W_{\leq n}^{(\tau,\sigma)} - \dim W_{<n}^{(\tau,\sigma)},
  \end{equation}
  where $W_{<n}^{(\tau,\sigma)}$ is the span generated by shells $S_0,\ldots,S_{n-1}$
  in proxy channel $(\tau,\sigma)$.
  The mean capacity $\bar{\Sigma}_n$ over sampled generic channels was shown to satisfy
  \begin{equation}
    \bar{\Sigma}_n \sim C\,n^{-\hat{\delta}_{\mathrm{cap}}},
    \qquad
    \hat{\delta}_{\mathrm{cap}} \approx 3.4 \text{ at } q=307,\; R^2=0.987.
    \label{eq:o11-result}
  \end{equation}
  All three measurability conditions (B1)--(B3) were satisfied at all tested primes.
  The proxy result~\eqref{eq:o11-result} provides the baseline against which
  $\hat{\delta}_{\mathrm{exact}}$ is compared in Section~\ref{sec:comparison}.

  \paragraph{Limitation of O11.}
    The converged value $\hat{\delta}_{\mathrm{cap}}\approx 3.4$ is a proxy-level estimate.
    Two interpretations were left open in O11: a proxy-level artefact due to the
    absence of~$\gamma$, or an insufficiently large~$q$.
    The first is addressed by the exact Weil projection at $q\in\{29,53,61\}$ in the present
    paper (Section~\ref{sec:comparison} resolves it in the negative: the central coordinate
    is exactly inert for this observable, not a proxy-level artefact).
    The second — insufficiently large~$q$ — is reframed rather than settled by finite-$q$
    computation: Section~\ref{sec:comparison} derives the intrinsic asymptotic exponent
    analytically, in closed form, making further finite-$q$ extension useful only as an
    empirical check, not as the route to the asymptotic value.

    This motivates the construction of an exact block-wise capacity in the full Weil
    representation, which we develop in Section~\ref{sec:exact-capacity}.
